
\documentclass{article}
\usepackage{colortbl}
\usepackage{makecell}
\usepackage{multirow}
\usepackage{supertabular}

\begin{document}

\newcounter{utterance}

\centering \large Interaction Transcript for game `cladder', experiment `full\_v1.5\_default', episode 4 with Qwen3.6{-}27B.
\vspace{24pt}

{ \footnotesize  \setcounter{utterance}{1}
\setlength{\tabcolsep}{0pt}
\begin{supertabular}{c@{$\;$}|p{.15\linewidth}@{}p{.15\linewidth}p{.15\linewidth}p{.15\linewidth}p{.15\linewidth}p{.15\linewidth}}
    \# & \multicolumn{2}{c}{Player} && \multicolumn{2}{c}{Game Master} \\
    \hline

    \theutterance \stepcounter{utterance}  
    & & & \multicolumn{4}{p{0.6\linewidth}}{
        \cellcolor[rgb]{0.9,0.9,0.9}{
            \makecell[{{p{\linewidth}}}]{
                \texttt{\tiny{[P1$\langle$GM]}}
                \texttt{You are an expert in causal inference. The following question is not a typical commonsense query, but rather a meticulously designed question created by a professor specializing in causal inference, intended to assess the students' mastery of the course content.} \\
\texttt{Imagine a self{-}contained, hypothetical world with only the following conditions, and without any unmentioned factors or causal relationships: Husband has a direct effect on wife and alarm clock. Wife has a direct effect on alarm clock. For husbands that don't set the alarm and wives that don't set the alarm, the probability of ringing alarm is 11\%. For husbands that don't set the alarm and wives that set the alarm, the probability of ringing alarm is 60\%. For husbands that set the alarm and wives that don't set the alarm, the probability of ringing alarm is 46\%. For husbands that set the alarm and wives that set the alarm, the probability of ringing alarm is 92\%. For husbands that don't set the alarm, the probability of alarm set by wife is 61\%. For husbands that set the alarm, the probability of alarm set by wife is 1\%. Does husband positively affect alarm clock through wife?} \\
\texttt{Start your answer with 'Yes' or 'No', followed by additional reasoning or evidence to support your explanation.} \\
            }
        }
    }
     \\ \\

    \theutterance \stepcounter{utterance}  
    & \multicolumn{4}{p{0.6\linewidth}}{
        \cellcolor[rgb]{0.561,0.737,0.561}{
            \makecell[{{p{\linewidth}}}]{
                \texttt{\tiny{[P1$\rangle$GM]}}
                \texttt{Yes} \\
\\ 
\texttt{To determine whether the husband positively affects the alarm clock through the wife, we need to calculate the **Indirect Effect** of the husband on the alarm clock via the wife. In causal inference, this is often referred to as the mediated effect.} \\
\\ 
\texttt{The causal structure is:} \\
\texttt{\$Husband (H) \textbackslash{}rightarrow Wife (W) \textbackslash{}rightarrow Alarm (A)\$} \\
\texttt{\$Husband (H) \textbackslash{}rightarrow Alarm (A)\$ (Direct effect)} \\
\\ 
\texttt{We are asked about the effect *through* the wife. This implies we are looking at the change in the probability of the alarm ringing caused by the husband's action, specifically mediated by the wife's action.} \\
\\ 
\texttt{Let's define the variables:} \\
\texttt{\$H=1\$: Husband sets alarm, \$H=0\$: Husband does not set alarm.} \\
\texttt{\$W=1\$: Wife sets alarm, \$W=0\$: Wife does not set alarm.} \\
\texttt{\$A=1\$: Alarm rings, \$A=0\$: Alarm does not ring.} \\
\\ 
\texttt{We are given the conditional probabilities for the alarm ringing (\$P(A=1 | H, W)\$):} \\
\texttt{1. \$P(A=1 | H=0, W=0) = 0.11\$} \\
\texttt{2. \$P(A=1 | H=0, W=1) = 0.60\$} \\
\texttt{3. \$P(A=1 | H=1, W=0) = 0.46\$} \\
\texttt{4} \\
            }
        }
    }
    & & \\ \\

    \theutterance \stepcounter{utterance}  
    & & & \multicolumn{2}{p{0.3\linewidth}}{
        \cellcolor[rgb]{0.9,0.9,0.9}{
            \makecell[{{p{\linewidth}}}]{
                \texttt{\tiny{[GM$|$GM]}}
                \texttt{Yes} \\
            }
        }
    }
    & & \\ \\

    \theutterance \stepcounter{utterance}  
    & & & \multicolumn{2}{p{0.3\linewidth}}{
        \cellcolor[rgb]{0.9,0.9,0.9}{
            \makecell[{{p{\linewidth}}}]{
                \texttt{\tiny{[GM$|$GM]}}
                \texttt{no} \\
            }
        }
    }
    & & \\ \\

    \theutterance \stepcounter{utterance}  
    & & & \multicolumn{2}{p{0.3\linewidth}}{
        \cellcolor[rgb]{0.9,0.9,0.9}{
            \makecell[{{p{\linewidth}}}]{
                \texttt{\tiny{[GM$|$GM]}}
                \texttt{game\_result = LOSE} \\
            }
        }
    }
    & & \\ \\

\end{supertabular}
}

\end{document}
